\section{Introduction}

\label{sec:introduction}

The Samsung Galaxy camera publishes an initial user-visible image, the \emph{draft image}, for each capture and later replaces it with the final image from post-processing.
Samsung refers to the processing stages that produce the draft image as the \emph{draft sequence}.
For certain shot modes, post-processing is deferred until the user switches away from the camera application.
The draft image then remains the user-visible result throughout the foreground session, making its visual quality more important.
In the draft sequence, Samsung has already deployed single-frame stages to narrow the draft--final visual gap and implemented lightweight multi-frame composition to reduce it further.

However, lightweight multi-frame composition has not been deployed on any device because its added processing can violate the \emph{Capture Timeout}, a fixed deadline for each draft image.
A violation crashes the application by design and blocks the release.
The camera must meet this deadline while retaining parallel capture, which shortens shot-to-shot intervals by accepting the next shutter input before the preceding draft image is produced.
Yet draft sequences still execute one at a time on a single draft worker, so later captures may have to wait for earlier draft sequences to finish.
Since each capture's timeout clock starts at request acceptance, added work in an earlier draft sequence can consume queued captures' timeout budgets and cause them to miss their deadlines.

Samsung already uses a static thermal guard to mitigate timeout risk from the deployed single-frame stages.
The guard omits optional stages when configuring the draft sequence at overheat level~4 or higher on a 0--6 scale.
However, in our motivating experiment with the guard bypassed, the capture index of the earliest timeout varied with resolution, memory pressure, and enabled draft stages, even at the same starting overheat level.
The overheat level reflects neither how much timeout budget a capture has left nor how much draft work remains.
Deploying multi-frame composition under parallel capture therefore requires runtime control based on remaining timeout budgets and draft workload.

Such runtime control is constrained: the draft pipeline can neither reject an accepted capture nor preempt a running stage.
Skipping optional stages reduces remaining processing time but cannot recover timeout budget already consumed by queueing.
Delaying the next capture opportunity limits further queue growth to preserve budget for subsequent captures.
These mechanisms correspond to workload control~\cite{approxdet-sensys,timecard-sosp} and arrival control~\cite{seda-sosp,breakwater-osdi}, respectively.
The two controls must be coordinated because decisions to execute or skip optional stages change the draft workload on which arrival control depends.

In this paper, we present the Budget-Aware Draft Controller, which coordinates these two controls to enable lightweight multi-frame composition within per-capture deadlines.
A major challenge is deciding how much time to reserve for draft work: too little risks a Capture Timeout, whereas too much needlessly suppresses optional processing or prolongs shot-to-shot intervals.
To address this, \emph{Remaining-Sequence Admission} uses recent stage overruns to calibrate an empirical upper duration estimate for the entire remaining draft sequence and admits an optional stage when this estimate fits within the capture's remaining timeout budget.
\emph{Capture-Availability Pacing} delays the next capture opportunity when projected draft workload exceeds the remaining timeout budget of the most recently enqueued capture.
Both use a shared duration predictor updated from measured durations in the current camera session, without per-device timing tables.

Across 28 Portrait-mode conditions on two commercial devices, the controller completed all 280 runs without Capture Timeout, whereas uncontrolled execution of all optional stages timed out in 225 of 280 runs.
Lightweight multi-frame composition executed on at least 4.5 of the first five captures per run on average in every condition, including those starting at or above the guard's cutoff.
With production-default 12MP capture under normal memory conditions, pacing did not delay the first five captures at starting overheat levels below 4.
Only the controller with both admission and pacing completed all captures on time in all three ablation conditions, while executing multi-frame composition 1.3--2.0 times as often as admission alone.
The controller is integrated into the product branch and is undergoing release validation for a forthcoming Galaxy flagship line.

\smallskip\noindent\textbf{Contributions.}
\begin{itemize}[leftmargin=1.2em]
  \item We characterize how serialized draft execution under parallel capture creates a cross-capture deadline problem for lightweight multi-frame composition, and why a static thermal guard does not reflect per-capture deadline safety.
  \item We present the Budget-Aware Draft Controller, which coordinates remaining-sequence admission with capture-availability pacing. To our knowledge, this is the first work to combine these controls in production camera software under a hard per-capture deadline.
  \item We evaluate the controller on two commercial devices, examining end-to-end effectiveness, component contribution, admission decision quality against naive approaches, and pacing under queue pressure.
\end{itemize}
